% 三线表骨架（booktabs + tabularx）
%
% 写之前确认两件事：
%   列的语义——每一列的单元格是不是同一种信息？
%   行的语义——每一行是不是在描述同一个对象？
% 只回答一个问题的小表通常用正文更清楚；完整论证不要塞进单元格。

\begin{table}[htbp]
\centering
\small
\caption{这张表回答什么问题}
\label{tab:example}
% V 列权重之和 = 本表 V 列的列数（此处 0.8 + 1.2 = 2）
\begin{tabularx}{\textwidth}{L{2.4cm}V{0.8}V{1.2}}
\toprule
\textbf{对象} & \textbf{字段一} & \textbf{字段二}\\
\midrule
第一个对象 & 取值 & 取值\\
第二个对象 & 取值 & 取值\\
\bottomrule
\end{tabularx}
\end{table}

% 图注回答「这是什么」，正文回答「它说明什么」；图、表、正文不重复播报同一过程。
% \begin{figure}[htbp]
% \centering
% \includegraphics[width=0.8\textwidth]{figure.pdf}
% \caption{图是什么；实现状态与适用范围}
% \label{fig:example}
% \end{figure}
